\documentclass[12pt]{article}
\usepackage[utf8]{inputenc}
\usepackage[T1]{fontenc}
\usepackage{hyperref}
\title{Warao Language Resources for K-12 Curriculum Development}
\author{Compiled by Claude Sonnet 5 (claude-sonnet-5) for Andr\'as Kornai}
\date{}
\begin{document}
\maketitle
\section{Introduction}
Warao (ISO 639-3 wba) is a language isolate of the Orinoco Delta in Venezuela and Guyana, with about 28,000 speakers. UNESCO considers it endangered. It is written in the Latin alphabet. Most speakers are bilingual in Spanish to varying degrees. One source says it is taught in some schools in Delta Amacuro State. The searches did not confirm any current school programmes.
\section{Primary Educational Resources}
No specific primers or textbooks were found in the searches. Omniglot gives a general description of the writing system.
\section{Linguistic and Reference Materials}
\begin{itemize}
\item Wikipedia overview: \url{https://en.wikipedia.org/wiki/Warao_language}
\item Andr\'es Romero-Figueroa, \emph{A Reference Grammar of Warao} (LINCOM, 2006), described in: \url{https://linguistlist.org/issues/17/3721/}
\item Leiden University source citing the Barral (1979) and Edwards (1980) dictionaries and the Vaquero (1965) and Romero-Figueroa (1997) grammar sketches: \url{https://scholarlypublications.universiteitleiden.nl/access/item%3A3576115/view}
\item Omniglot, writing system: \url{https://www.omniglot.com/writing/warao.htm}
\end{itemize}
\section{Religious and Community Resources}
The searches turned up no scripture or community organisation pages. The Barral missionary history is mentioned in the Wikipedia article.
\section{Digital and Audio Resources}
No audio or software resources were found in the searches. A Brazilian journal article appeared in the results: \url{https://periodicos.ufpb.br/index.php/enid/article/download/74754/41774/233851}
\section{Curriculum Development Context}
Warao has a reference grammar and older dictionaries, which give a basis for materials. Current bilingual education policy and school materials need further investigation.
\end{document}